\section{Alcance del Proyecto}

Este proyecto se centra en construir y evaluar un modelo de aprendizaje automático
capaz de estimar el nivel de riesgo de contratos públicos de infraestructura justo
cuando se adjudican, antes de que comience la ejecución. La materia prima son los
registros abiertos de SECOP~II, particularmente los contratos que se financian con
recursos del Sistema General de Regalías~(SGR).

\subsection{Entregables Funcionales}

Al cierre del proyecto se contará con los siguientes productos concretos:

\begin{itemize}
    \item Un \textit{pipeline} de extracción y preparación de datos que cubra desde
    la descarga de registros de SECOP~II hasta la generación de un conjunto de datos
    listo para modelar. Esto incluye manejo de nulos, corrección de inconsistencias
    territoriales y la construcción indirecta de la variable de riesgo a través de
    indicadores \textit{proxy}.

    \item Un módulo de ingeniería de características con las variables que la
    literatura identifica como más informativas: porcentaje de anticipo, modalidad
    de contratación, número de proponentes y localización del contrato.

    \item Un módulo de entrenamiento que compare tres clasificadores ---Regresión
    Logística, Random Forest y XGBoost--- ajustados con búsqueda aleatoria de
    hiperparámetros (\texttt{RandomizedSearchCV}) para mantenerse dentro de los
    límites computacionales del proyecto.

    \item Un módulo de evaluación con validación cruzada estratificada ($k = 5$) y
    reporte de AUC-ROC, F1-score, precisión y \textit{recall}. Se añadirá un análisis
    de importancia de variables para saber qué factores pesan más en la predicción.

    \item El modelo seleccionado, serializado y documentado, junto con el código
    reproducible del experimento completo.

    \item Un informe técnico que compare los tres algoritmos, discuta qué variables
    resultaron más relevantes y proponga cómo podría escalar este trabajo hacia un
    Sistema de Alertas Tempranas operativo.
\end{itemize}

\subsection{Área de Estudio}

El foco está puesto en contratos de infraestructura del SGR registrados en SECOP~II.
Se eligió ese universo por dos razones prácticas: los montos involucrados son
considerables y la plataforma ofrece datos estructurados de acceso público, lo que
hace viable el proyecto sin necesidad de convenios institucionales.

\subsection{Frontera del Proyecto}

Para no exceder la capacidad del equipo en el tiempo disponible, quedan fuera del
alcance las siguientes actividades:

\begin{itemize}
    \item Desplegar el modelo en producción o integrarlo con los sistemas de alguna
    entidad de control. El trabajo termina en la validación académica del modelo.

    \item Desarrollar cualquier tipo de interfaz de usuario ---web, móvil o de
    escritorio---. Los resultados se presentan en el informe y el repositorio de
    código.

    \item Ejecutar búsquedas exhaustivas de hiperparámetros (Grid Search completo),
    dado el costo computacional que implica sobre XGBoost con espacios de parámetros
    amplios.

    \item Detectar esquemas de corrupción que no dejen huella en los registros
    públicos: acuerdos verbales, testaferros no vinculados o pagos fuera del sistema.
    Eso requiere fuentes de inteligencia financiera que están fuera del alcance de
    este trabajo.
\end{itemize}

